\chapter{Sesión 6c: Publicaciones y contadores de no leídos en Flutter}
\label{ch:sesion6c}

\section{Objetivo de la sesión}

Completar el cliente móvil con:

\begin{itemize}[leftmargin=*]
  \item Visualización de las publicaciones dentro del canal.
  \item Publicación de avisos desde la app.
  \item Consulta y despliegue de contadores de no leídos por canal.
  \item Marcado automático como leído al abrir un canal.
  \item Badge numérico en la pestaña ``Suscritos'' con el total de no leídos.
\end{itemize}

\section{Contexto}

La Sesión~\ref{ch:sesion6b} dejó el listado, búsqueda, detalle y creación de canales funcional. Sin embargo, el detalle del canal solo mostraba un placeholder en lugar de las publicaciones. El backend (Sesiones~\ref{ch:sesion4} y~\ref{ch:sesion5}) ya exponía endpoints de posts y contadores, así que el trabajo era exclusivamente del cliente.

\section{Desarrollo}

\subsection{Archivos nuevos}

\begin{lstlisting}[language=bash]
mobile/lib/state/
|- post_provider.dart           (NUEVO)
`- unread_provider.dart         (NUEVO)
\end{lstlisting}

\subsection{Archivos modificados}

\begin{itemize}[leftmargin=*]
  \item \texttt{lib/main.dart}: se registraron \texttt{PostProvider} y \texttt{UnreadProvider} en el \texttt{MultiProvider}.
  \item \texttt{lib/screens/channel\_detail\_screen.dart}: se reemplazo el placeholder por la lista de publicaciones, la caja para publicar y la logica de marcado como leido.
  \item \texttt{lib/screens/home\_shell.dart}: se agrego el badge numerico en la pestana ``Suscritos''.
\end{itemize}

\subsection{Decisiones tecnicas}

\begin{itemize}[leftmargin=*]
  \item \textbf{Un \texttt{PostProvider} que cachea por canal.} Mantiene un \texttt{Map<String, List<Post>>} indexado por \texttt{channelId}. Esto permite cambiar entre canales sin rehacer peticiones.
  \item \textbf{Carga optimista al publicar.} Cuando el usuario publica, el post se inserta inmediatamente al inicio de la lista local sin esperar a recargar todo el canal.
  \item \textbf{Borrado local tras exito del backend.} El \texttt{delete} solo modifica el estado local despues de que el servidor responde 204.
  \item \textbf{Contadores en Redis como fuente de verdad.} El \texttt{UnreadProvider} solo consulta \texttt{/api/me/unread} y expone el \texttt{Map<String, int>} resultante.
  \item \textbf{Marcado como leido al entrar.} Al cargar el detalle del canal, se llama automaticamente a \texttt{markRead(channelId)}.
\end{itemize}

\subsection{Componentes clave}

\subsubsection{\texttt{PostProvider}}

\begin{lstlisting}[language=Java]
class PostProvider extends ChangeNotifier {
  final Map<String, List<Post>> _byChannel = {};
  final Map<String, bool> _loading = {};

  List<Post> postsOf(String channelId);
  bool isLoading(String channelId);

  Future<void> load(String channelId, {bool refresh = false});
  Future<Post> create(String channelId, String content);
  Future<void> delete(String channelId, String postId);
}
\end{lstlisting}

\subsubsection{\texttt{UnreadProvider}}

\begin{lstlisting}[language=Java]
class UnreadProvider extends ChangeNotifier {
  Map<String, int> _counts = {};

  int get total => _counts.values.fold(0, (a, b) => a + b);
  int countFor(String channelId) => _counts[channelId] ?? 0;

  Future<void> refresh();
  Future<void> markRead(String channelId);
  void clear();
}
\end{lstlisting}

\subsubsection{Pantalla de detalle del canal}

Estructura visual:

\begin{itemize}[leftmargin=*]
  \item Encabezado: nombre, chip de categoria, badge de rol del owner (se anade en la Fase~E), descripcion, autor, contador de suscriptores y boton Suscribirse/Desuscribirse.
  \item Caja de publicacion: \texttt{TextField} multilinea + boton de envio.
  \item Lista de publicaciones con \texttt{RefreshIndicator} y paginacion inicial de 50 elementos.
\end{itemize}

\subsection{Cambios a archivos previos}

\begin{itemize}[leftmargin=*]
  \item \texttt{main.dart}: \texttt{MultiProvider} ampliado con dos providers.
  \item \texttt{channel\_detail\_screen.dart}: reemplazo completo.
  \item \texttt{home\_shell.dart}: badge en la pestana.
\end{itemize}

\subsection{Problemas encontrados y solucion}

\subsubsection{Locale de fechas no inicializado}

Al abrir el detalle del canal, la app mostraba:

\begin{lstlisting}[language=bash]
LocaleDataException: Locale data has not been initialized,
call initializeDateFormatting(<locale>).
\end{lstlisting}

\textbf{Solucion:} inicializar el locale en \texttt{main()}:

\begin{lstlisting}[language=Java]
Future<void> main() async {
  WidgetsFlutterBinding.ensureInitialized();
  await initializeDateFormatting('es', null);
  runApp(const NotiuamApp());
}
\end{lstlisting}

\section{Verificacion}

Con el backend corriendo en \texttt{localhost:8081} y una sesion iniciada:

\begin{enumerate}[leftmargin=*]
  \item Entrar a cualquier canal.
  \item Ver el listado de publicaciones anteriores.
  \item Publicar un aviso nuevo desde la caja de texto.
  \item Confirmar que aparece al instante en la lista.
  \item Abrir el canal desde el usuario destinatario y verificar que el badge de no leidos se limpia.
\end{enumerate}

\section{Resultado}

\begin{exito}[Estado al cierre]
\begin{itemize}[leftmargin=*]
  \item Publicaciones visibles y creables desde la app.
  \item Contadores de no leidos consumidos desde Redis.
  \item Badge numerico en la barra inferior.
  \item Marcado automatico como leido al abrir el canal.
  \item Base para la moderacion directa (Fase~D5) y para push real (Fase~F).
\end{itemize}
\end{exito}